%1/2in left/right margins, 1in top, 1/2 bottom
\documentclass[11pt]{article}
\hbadness=10000
\tolerance=10000
\textheight 9.7in
\textwidth 7.4in
\oddsidemargin -.5in
\topmargin -1in
\headsep 0in
\headheight 0in
\begin{document}
\hrule
\medskip
\centerline{\textbf{STATISTICS 801}}\hfill \textbf{Fall 2016}
\centerline{\textbf{Extra problems for paired data and inference about the population variance}}
\smallskip
\hrule
\smallskip
\begin{enumerate}
\item Consider the following study in which standing and supine systolic blood pressures were compared. To eliminate the variation in blood pressure among subjects affecting the accuracy of the conclusions drawn, the blood pressures were measured on twelve subjects, each measured in both positions. It is therefore, a paired samples design. \\[-.2in]
\begin{center}
\begin{tabular}{l|@{\extracolsep{2mm}}rrrrrrrrrrrr}\hline
Subject & 1 & 2 & 3 & 4 & 5 & 6 & 7 & 8 & 9 & 10 & 11 & 12  \\ \hline
Standing & 132 & 146 & 135 & 141  & 139 & 162 & 128 & 137 & 145 & 151 & 131 & 143 \\
Supine   & 136 & 145 & 140 & 147  & 142 & 160 & 137 & 136 & 149 & 158 & 120 & 150 \\ \hline
\multicolumn{1}{c|}{$d$}  \\ \hline
\multicolumn{1}{c|}{$d^2$} \\ \hline
\end{tabular}
\end{center}


\texttt{Note:} Do all computations for parts (b),(c), and (d) using hand calculation. 
Use the output from your analysis using the data
 to provide answers to parts (a),(e) and (f). \\[-.2in] 
\begin{enumerate}

\item Do the Normal probability plot and the boxplot from the output indicate that a
paired t-test could be used? Explain. Obtain a plot of the data that shows the two variables are positively correlated.
What does this say about this this design compaired to a two-independent samples study? \\[1.5in]

\item
Compute the paired t-statistic to test the research hypotheses that 
there is no difference in mean blood pressure in the two positions. Use $\alpha = .05$.
State your hypotheses clearly. Define the symbols you use carefully.\\[1.5in]
\item Explain what is meant by the p-value in words by defining it for the above test.  
Show how to find bounds for the p-value of the above test using the t-table. 

\newpage
\item
Use a 95\% confidence interval to estimate the difference in mean blood pressure in the two positions. 
Test the hypothesis in part (b) using this
interval stating the $\alpha$ level used for this test clearly.
State your conclusion.\\[2in]

\item Extract the t-statistic from the your output for testing the
hypothesis you stated in part (b). 
State the p-value associated with this test extracted from the output.\\[.8in]

\item  Extract the  95\% confidence interval for the  difference in mean blood pressure in the two positions  from the output. 
\\[.8in]

\end{enumerate}

\item The times of first sprinkler activation for a series of tests
with fire prevention sprinkler systems using an aqueous film-forming
foam were (in sec)
$ 27, 32, 22, 27, 23, 35, 30, 33, 24, 27, 21, 22, 24 $.\\
The system has been manufactured under the design specification 
that the true standard deviation of the activation time will not  exceed 5 sec under 
normal conditions. Given $\sum{y}=347$ and $\sum{y^2}=9515$\\
{\tt Note:}  Use hand calculation for parts (b) and (c) and show
work. Use a computer analysis of the data for answering parts (a),(d) and (e).


\begin{enumerate}
\item Do the box plot or the normal probability plot indicate any violation of the conditions
underlying the use of the chi-square procedures for constructing
confidence intervals or testing hypotheses about $\sigma$?
\newpage
\item  Is there sufficient evidence in the data to determine that the standard deviation of the sprinkler activation time meets the design specification? State appropriate hypotheses and perform a test with $\alpha=.05$.\\[2in]

\item Estimate the standard deviation  of the sprinkler activation time
using a 95\% confidence interval.\\[2in]


\item Extract the statistic from the your output for testing the
hypothesis you stated in part (b).
State the p-value associated with this test extracted from the  output.\\[1in]

\item  Extract the  95\% confidence interval for the standard deviation 
of the sprinkler activation time from the output. 
\end{enumerate}






\newpage
\item Twenty plants of the same variety were grown from cuttings. Ten plants, selected randomly,
were grown in soil type A and the others in soil type B. In all respects, other than
soil type, the cuttings were cared for in the same way. At the end of the experiment the
plants were uprooted and dried. The values in the table below are their dry weights (in grams).
Note that one cutting (in soil type A) died at an early stage in the experiment.
\begin{center}
\begin{tabular}{c|rrrrrrrrrr|rr}
\hline 
 \multicolumn{1}{c|}{\textit{Soil Type}} & \multicolumn{10}{c|}{\textit{Dry Weight(in grams)}} & 
     \multicolumn{1}{c}{$\sum{y}$} &
     \multicolumn{1}{c}{$\sum{y^2}$}  \\ \hline
A & 25.5 &  22.3 &  24.7 &  28.1 &  26.5 &  19.0 &  23.0 &  25.3 &  29.2 & &223.6  & 5632.22     \\
B & 31.8 &  30.3 &  26.4 &  24.2 &  27.8 &  29.1 &  25.5 &  28.9 &  26.9 &  29.7 & 280.6 & 7922.74 \\ \hline
\end{tabular}
\end{center}
\begin{enumerate}
\item Use the boxplot and the normal probabilty plot produced in the computer output
to comment on whether the assumption of normality of the two populations
is reasonable. Do these plots also support the assumption that
the population variances are the same for the two populations? 
Explain\\[1in]


\item Compute the sample variances $s_A^2$ and $s_B^2$, respectively, of the two samples using the above summary statistics.\\[.8in]
\item Compute a statistic to test the research hypothesis that the two 
population variances of dry weights for the two soil types are different using 
$\alpha=.05$ by hand. State your hypotheses, the rejection region and 
your conclusion.\\[1.8in]


\item  Find the statistic for testing the hypothesis you stated in part (c) in the computer output. State the p-value associated with this test extracted from the output here. \\[1in]
 \end{enumerate}

\end{enumerate}
\end{document}


